\section{第41届复赛第6题：太阳引力场中的轨道偏折与星光弯曲}
\subsection{题干}
题目先在牛顿动力学框架下讨论圆锥曲线轨道 $r(\theta)=p/(1+e\cos\theta)$，再分析以瞄准距离 $b$ 进入太阳系的陨石和掠日星光的偏折；最后借助弱场有效折射率和费马原理重新计算星光偏折。
\begin{figure}[htbp]\centering\begin{tikzpicture}[>=Stealth,scale=.88]
  \draw[densely dotted,thick] (-4.2,1.15)--(4.5,1.15); \draw[densely dotted,thick] (-4.2,-1.15)--(.65,-1.15);
  \fill (.65,-1.15) circle(.09); \node[below] at (.65,-1.25) {$M_\odot$};
  \draw[very thick] (-4.0,1.15).. controls (-.2,1.18) and (1.25,.92)..(2.1,.25)..controls(3.0,-.15) and (3.8,-.62)..(4.4,-.98);
  \draw[->,very thick] (-4.15,1.15)--(-3.25,1.15) node[midway,above] {$\boldsymbol v_0$};
  \draw[<->] (-2.8,-1.15)--(-2.8,1.15) node[midway,left] {$b$};
  \draw[densely dotted,thick] (1.9,1.18)--(4.35,1.18);
  \draw (2.05,1.15) arc[start angle=0,end angle=-26,radius=1.25]; \node[right] at (2.85,.88) {$\Delta\varphi$};
  \node[below] at (.2,-1.9) {原题图6a：掠过太阳附近的入射轨迹与偏折角};
\end{tikzpicture}
\caption{第41届复赛第6题原题图6a的TikZ重绘。}\end{figure}
\subsection{参考解答整理}
\paragraph{（1）双曲轨道的角度范围。}
对 $e>1$，轨道可取范围由 $\cos\theta>-1/e$ 给出，故
\[\theta_{\max}=\arccos(-1/e),\qquad \frac\pi2<\theta_{\max}<\pi.\]
\paragraph{（2）用能量和角动量表示轨道参数。}
以总能量 $E$ 和角动量 $L$ 表示的轨道参数为
\[p=\frac{L^2}{GMm^2},\qquad e=\sqrt{1+\frac{2EL^2}{G^2M^2m^3}}.\]
\paragraph{（3）牛顿框架下的陨石和星光偏折。}
对速度 $v_0$、瞄准距离 $b$ 的入射陨石，在 $v_0^2b/(GM_\odot)\gg1$ 条件下，偏折角为
\[\boxed{\Delta\varphi_{\mathrm{meteor}}=\frac{2GM_\odot}{v_0^2b}}.\]
令 $v_0=c$、$b=R_\odot$，牛顿框架下的星光偏折是
\[\Delta\varphi_{\mathrm{Newton}}=\frac{2GM_\odot}{c^2R_\odot}.\]
\paragraph{（4）弱场广义相对论框架下的星光偏折。}
弱场广义相对论以 $n_{\mathrm{eff}}(r)=1-2\phi(r)/c^2$、$\phi(r)=-GM_\odot/r$ 描述光程。由费马原理导出的守恒式与无穷远边界条件给出最终偏折
\[\boxed{\Delta\varphi_{\mathrm{GR}}=\frac{4GM_\odot}{c^2R_\odot}=2\Delta\varphi_{\mathrm{Newton}}}.\]
\subsection{转写依据}
题干依据 \texttt{CPHO41-Q06\_transcription\_v1.md}；答案公式依据答案第13至16页和 \texttt{CPHO41-Q06\_formula\_check\_v1.csv}。图6a由原题图重绘，不嵌入原始页面截图。
